% =============================================================================
% APPENDIX: DOMAIN-AIMARK: AI Marketing Transformation Axioms
% =============================================================================
% Module: MOD-UBS-AIMARK-001
% Version: 1.0 (February 2026)
% Status: Draft
%
% Complete axiom system for the UBS AI Marketing Transformation Model.
% Formalizes 4 modules with 20 axioms, 12 definitions, 8 propositions.
% Monte Carlo validated (10,000 draws, all propositions confirmed).
%
% =============================================================================

\chapter{AIMARK: DOMAIN-AIMARK: AI Marketing Transformation --- Formal Axiom System}
\label{app:aimark}

\begin{abstract}
This appendix provides the complete formal axiom system for the AI Marketing
Transformation Model (MOD-UBS-AIMARK-001). The model integrates four modules
--- AI Value Chain, Customer Scenarios, Organization Transformation, and
Strategic Roadmap --- into a unified framework for behavioral analysis of
AI-driven marketing transformation in wealth management. We present 20 axioms
(AIM-1 through AIM-20), 12 definitions, and 8 propositions, all validated
via Monte Carlo simulation with 10,000 draws.
\end{abstract}

%%%%%%%%%%%%%%%%%%%%%%%%%%%%%%%%%%%%%%%%%%%%%%%%%%%%%%%%%%%%%%%%%%%%%%%%%%%%%%%
% QUICK REFERENCE
%%%%%%%%%%%%%%%%%%%%%%%%%%%%%%%%%%%%%%%%%%%%%%%%%%%%%%%%%%%%%%%%%%%%%%%%%%%%%%%

\begin{tcolorbox}[colback=blue!5!white, colframe=blue!75!black,
    title=Quick Reference: AI Marketing Transformation Axioms]

\textbf{Core Question:} How does AI transform marketing value creation,
customer behavior, organizational capability, and strategic optionality
in wealth management?

\textbf{Key Formula:}
\begin{equation}
U_{\text{total}} = \sum_{m} w_m(\mathcal{S}) \cdot V_m + \sum_{i < j} \gamma_{ij} \cdot \sqrt{V_i \cdot V_j}
\label{eq:aim-integrated}
\end{equation}

\textbf{Axiom Count:} 20 axioms (AIM-1 through AIM-20)

\textbf{Modules:}
\begin{itemize}[nosep]
\item \textbf{M1:} AI Value Chain (AIM-1 through AIM-5)
\item \textbf{M2:} Customer Scenarios (AIM-6 through AIM-10)
\item \textbf{M3:} Organization Transformation (AIM-11 through AIM-15)
\item \textbf{M4:} Strategic Roadmap (AIM-16 through AIM-20)
\end{itemize}

\textbf{Cross-References:} Appendices B (HOW), V (WHEN), BBB (WHERE),
AV (READY), AW (STAGE), IE (EIT)

\end{tcolorbox}

%%%%%%%%%%%%%%%%%%%%%%%%%%%%%%%%%%%%%%%%%%%%%%%%%%%%%%%%%%%%%%%%%%%%%%%%%%%%%%%
% SCOPE BOX
%%%%%%%%%%%%%%%%%%%%%%%%%%%%%%%%%%%%%%%%%%%%%%%%%%%%%%%%%%%%%%%%%%%%%%%%%%%%%%%

\begin{tcolorbox}[
    colback=orange!5!white,
    colframe=orange!75!black,
    title={\textbf{Appendix Scope: Ziel / In-Scope / Out-of-Scope / Constraints}},
    fonttitle=\bfseries
]

\textbf{Ziel (Objective):}
\begin{itemize}[nosep]
\item Provide a complete, self-contained axiom system for AI-driven marketing transformation in wealth management
\end{itemize}

\textbf{In-Scope:}
\begin{itemize}[nosep]
\item Formal definitions of all model variables and parameters
\item 20 axioms with mathematical statements, interpretations, and testable predictions
\item 8 propositions derived from the axiom system with Monte Carlo validation
\item Cross-module complementarity formalization with EXC-5 justification
\end{itemize}

\textbf{Out-of-Scope:}
\begin{itemize}[nosep]
\item Empirical calibration with primary data $\rightarrow$ Appendix BBB (WHERE)
\item Generic complementarity theory $\rightarrow$ Appendix B (HOW)
\item Skill formation technology proofs $\rightarrow$ Appendix SF (Dynamic Complementarity)
\end{itemize}

\textbf{Constraints:}
\begin{itemize}[nosep]
\item Parameters are LLMMC priors until empirical validation (Q3 2026)
\item Scenario probabilities are subjective Bayesian (Dirichlet prior)
\item Cross-module $\gamma$ values are analogical (from PAR-COMP-001 family)
\end{itemize}

\textbf{Lieferobjekte (Deliverables):}
\begin{itemize}[nosep]
\item Complete axiom system (20 axioms)
\item Proposition set with Monte Carlo validation (8 propositions)
\item Formal notation table for all symbols
\end{itemize}

\end{tcolorbox}

%%%%%%%%%%%%%%%%%%%%%%%%%%%%%%%%%%%%%%%%%%%%%%%%%%%%%%%%%%%%%%%%%%%%%%%%%%%%%%%
% CHAPTER LINKAGE
%%%%%%%%%%%%%%%%%%%%%%%%%%%%%%%%%%%%%%%%%%%%%%%%%%%%%%%%%%%%%%%%%%%%%%%%%%%%%%%

\begin{tcolorbox}[colback=purple!5!white, colframe=purple!75!black,
    title=Chapter Linkage]

\textbf{Primary Chapter:} Chapter 14: Behavioral Strategy

\textbf{Secondary Chapters:}
\begin{itemize}[nosep]
\item Chapter 5: HOW --- Complementarity (\S5.3: Cross-domain $\gamma$)
\item Chapter 9: WHEN --- Context (\S9.2: $\Psi$ in financial services)
\item Chapter 12: READY --- Skill readiness (\S12.4: Dynamic complementarity)
\item Chapter 17: Intervention Design (\S17.5: No-regret scoring)
\end{itemize}

\textbf{Reading Guide:}
\begin{itemize}[nosep]
\item For conceptual overview: Read Monte Carlo report first (EBF-OUT-FIN-001)
\item For formal details: This appendix provides complete specification
\item For parameter sources: See Appendix BBB (WHERE)
\end{itemize}

\end{tcolorbox}


%%%%%%%%%%%%%%%%%%%%%%%%%%%%%%%%%%%%%%%%%%%%%%%%%%%%%%%%%%%%%%%%%%%%%%%%%%%%%%%
\section{Notation and Preliminaries}
\label{sec:aim-notation}
%%%%%%%%%%%%%%%%%%%%%%%%%%%%%%%%%%%%%%%%%%%%%%%%%%%%%%%%%%%%%%%%%%%%%%%%%%%%%%%

\begin{table}[htbp]
\centering
\caption{Complete Symbol Table for MOD-UBS-AIMARK-001}
\label{tab:aim-symbols}
\renewcommand{\arraystretch}{1.3}
\begin{tabular}{@{}cll@{}}
\toprule
\textbf{Symbol} & \textbf{Name} & \textbf{Domain / Scale} \\
\midrule
\multicolumn{3}{l}{\textit{Module 1: AI Value Chain}} \\
$\mathcal{F}$ & Set of marketing functions & $\{F_1, \ldots, F_7\}$ \\
$\alpha_f$ & AI potential of function $f$ & $[0,1]$ \\
$\varphi_f$ & Implementation feasibility & $[0,1]$ \\
$\delta_f$ & Differentiation premium & $[0,1]$ \\
$w_f$ & Strategic weight of function $f$ & $[0,1], \sum w_f = 1$ \\
$V_f$ & AI value of function $f$ & $\mathbb{R}_{\geq 0}$ \\
\midrule
\multicolumn{3}{l}{\textit{Module 2: Customer Scenarios}} \\
$\mathcal{S}$ & Scenario space & $\{S_1, S_2, S_3\}$ \\
$P(S_k)$ & Scenario probability & $[0,1], \sum P(S_k) = 1$ \\
$\boldsymbol{\alpha}$ & Dirichlet concentration & $\mathbb{R}_{>0}^3$ \\
$\mathcal{G}$ & Set of customer segments & $\{g_1, \ldots, g_4\}$ \\
$\tau_g$ & Trust parameter of segment $g$ & $[0,1]$ \\
$\lambda_g$ & Loss aversion of segment $g$ & $[1, \infty)$ \\
$\beta_g$ & Present bias of segment $g$ & $(0,1]$ \\
$\sigma_g$ & Switching cost of segment $g$ & $[0,1]$ \\
$\eta_g$ & AI adoption rate & $[0,1]$ \\
$\phi_g$ & Fear of asking (Frageangst) & $[0,1]$ \\
$R_g$ & Retention probability & $[0,1]$ \\
$n_g$ & Segment size (clients) & $\mathbb{N}$ \\
\midrule
\multicolumn{3}{l}{\textit{Module 3: Organization Transformation}} \\
$K_{\text{AI}}(t)$ & AI skill capital at time $t$ & $\mathbb{R}_{\geq 0}$ \\
$K_{\text{H}}(t)$ & Human skill capital at time $t$ & $\mathbb{R}_{\geq 0}$ \\
$I(t)$ & Training investment at time $t$ & $\mathbb{R}_{\geq 0}$ \\
$\theta(t)$ & Productivity of skill investment & $(0,1)$ \\
$\xi$ & DC premium (early vs.\ late) & $\mathbb{R}_{>0}$ \\
\midrule
\multicolumn{3}{l}{\textit{Module 4: Strategic Roadmap}} \\
$\rho_j$ & Scenario robustness of move $j$ & $[0,1]$ \\
$V_j^{\min}$ & Minimum value across scenarios & $\mathbb{R}_{\geq 0}$ \\
$O_j$ & Option value of move $j$ & $[0,1]$ \\
$r_j$ & Reversibility of move $j$ & $[0,1]$ \\
$L_j$ & Learning value of move $j$ & $[0,1]$ \\
$\text{NR}_j$ & No-regret score of move $j$ & $\mathbb{R}_{\geq 0}$ \\
$T_i(t)$ & Trigger signal $i$ at time $t$ & $\mathbb{R}$ \\
$\bar{T}_i^{(k)}$ & Threshold for trigger $i$, scenario $k$ & $\mathbb{R}$ \\
\midrule
\multicolumn{3}{l}{\textit{Cross-Module}} \\
$\gamma_{ij}$ & Complementarity between modules $i,j$ & $[-1,1]$ \\
$V_m$ & Value of module $m$ & $\mathbb{R}_{\geq 0}$ \\
$w_m(\mathcal{S})$ & Module weight (scenario-dependent) & $[0,1], \sum w_m = 1$ \\
$U_{\text{total}}$ & Integrated portfolio value & $\mathbb{R}$ \\
\bottomrule
\end{tabular}
\end{table}


%%%%%%%%%%%%%%%%%%%%%%%%%%%%%%%%%%%%%%%%%%%%%%%%%%%%%%%%%%%%%%%%%%%%%%%%%%%%%%%
\section{Module 1: AI Value Chain Axioms}
\label{sec:aim-m1}
%%%%%%%%%%%%%%%%%%%%%%%%%%%%%%%%%%%%%%%%%%%%%%%%%%%%%%%%%%%%%%%%%%%%%%%%%%%%%%%

\begin{tcolorbox}[colback=gray!5!white, colframe=gray!75!black,
    title=Axiom AIM-1: Multiplicative Value Decomposition]
\label{ax:aim-1}

\textbf{Statement:}
The AI value of a marketing function $f \in \mathcal{F}$ decomposes multiplicatively:
\begin{equation}
V_f = w_f \cdot \alpha_f \cdot \varphi_f \cdot (1 + \delta_f)
\label{eq:aim-value-decomposition}
\end{equation}

\textbf{Interpretation:} Value requires all three components simultaneously.
High AI potential ($\alpha_f$) is worthless without feasibility ($\varphi_f$);
both are necessary but not sufficient without strategic weight ($w_f$).
Differentiation ($\delta_f$) provides a multiplicative premium.

\textbf{EXC-5 Justification:} V1 (Veto) --- $\alpha_f = 0 \Rightarrow V_f = 0$
regardless of feasibility. P1 (Unit consistency) --- dimensions require
multiplication.

\textbf{Testable Prediction:} Functions with $\alpha_f > 0.7$ but $\varphi_f < 0.3$
will show near-zero realized value despite high AI potential.
\end{tcolorbox}


\begin{tcolorbox}[colback=gray!5!white, colframe=gray!75!black,
    title=Axiom AIM-2: Exhaustive Function Space]
\label{ax:aim-2}

\textbf{Statement:}
The marketing value chain is partitioned into $|\mathcal{F}| = 7$ mutually
exclusive and collectively exhaustive functions:
\begin{equation}
\mathcal{F} = \{F_1^{\text{Content}}, F_2^{\text{Campaign}}, F_3^{\text{Intelligence}},
F_4^{\text{Performance}}, F_5^{\text{Discovery}}, F_6^{\text{Competitive}}, F_7^{\text{Crisis}}\}
\label{eq:aim-function-space}
\end{equation}
with $\sum_{f \in \mathcal{F}} w_f = 1$ and $w_f > 0 \;\; \forall f$.

\textbf{Interpretation:} Every marketing activity maps to exactly one function.
No function has zero strategic weight (all are necessary for a complete
marketing organization).

\textbf{Testable Prediction:} Any new marketing capability (e.g., AI-generated
video) can be classified into exactly one $F_f$.
\end{tcolorbox}


\begin{tcolorbox}[colback=gray!5!white, colframe=gray!75!black,
    title=Axiom AIM-3: AI Potential Ordering]
\label{ax:aim-3}

\textbf{Statement:}
There exists a stable partial ordering of AI potentials:
\begin{equation}
\alpha_{F_3} > \alpha_{F_5} > \alpha_{F_1} > \alpha_{F_4}
\label{eq:aim-ordering}
\end{equation}
This ordering is robust under parameter uncertainty:
$P(\alpha_{F_3} > \alpha_{F_5}) \geq 0.85$.

\textbf{Interpretation:} Customer Intelligence ($F_3$) and AI-Mediated Discovery
($F_5$) are structurally more amenable to AI transformation than creative ($F_1$)
or performance ($F_4$) functions, because they operate on structured data
with clear optimization objectives.

\textbf{MC Validation:} $F_3$ ranked \#1 in 99.8\% of 10,000 draws;
$F_5$ ranked \#2 in 85.5\%.
\end{tcolorbox}


\begin{tcolorbox}[colback=gray!5!white, colframe=gray!75!black,
    title=Axiom AIM-4: Differentiation Asymmetry]
\label{ax:aim-4}

\textbf{Statement:}
The differentiation premium is asymmetrically distributed:
\begin{equation}
\delta_{F_5} > \delta_{F_3} > \delta_{F_1} \gg \delta_{F_2} \approx \delta_{F_4}
\label{eq:aim-diff-asymmetry}
\end{equation}

\textbf{Interpretation:} AI-Mediated Discovery ($F_5$) offers the highest
differentiation because it creates a \emph{new channel} (AI assistants recommending
UBS), while Performance Marketing ($F_4$) is commoditized (all competitors
optimize the same platforms).

\textbf{Theory Support:} MS-AI-001 (AI Consumer Experience) --- communicating
capability generates higher differentiation than operating capability.
\end{tcolorbox}


\begin{tcolorbox}[colback=gray!5!white, colframe=gray!75!black,
    title=Axiom AIM-5: Total Value Additivity]
\label{ax:aim-5}

\textbf{Statement:}
Total Module 1 value is additive across functions:
\begin{equation}
V_{M1} = \sum_{f \in \mathcal{F}} V_f = \sum_{f \in \mathcal{F}} w_f \cdot \alpha_f \cdot \varphi_f \cdot (1 + \delta_f)
\label{eq:aim-m1-total}
\end{equation}

\textbf{EXC-1 Justification:} Additivity is the default (EXC-1).
Intra-module function complementarities are captured at the cross-module
level via $\gamma_{M1, M3}$. No veto relationship between functions
(V1 not satisfied: $V_{F_3} = 0 \not\Rightarrow V_{F_1} = 0$).

\textbf{MC Validation:} $V_{M1} = 0.90 \; [0.75, 1.07]$ (95\% CI).
\end{tcolorbox}


%%%%%%%%%%%%%%%%%%%%%%%%%%%%%%%%%%%%%%%%%%%%%%%%%%%%%%%%%%%%%%%%%%%%%%%%%%%%%%%
\section{Module 2: Customer Scenario Axioms}
\label{sec:aim-m2}
%%%%%%%%%%%%%%%%%%%%%%%%%%%%%%%%%%%%%%%%%%%%%%%%%%%%%%%%%%%%%%%%%%%%%%%%%%%%%%%

\begin{tcolorbox}[colback=gray!5!white, colframe=gray!75!black,
    title=Axiom AIM-6: Scenario Exhaustiveness and Dirichlet Prior]
\label{ax:aim-6}

\textbf{Statement:}
The scenario space $\mathcal{S} = \{S_1, S_2, S_3\}$ is collectively exhaustive
and mutually exclusive. Scenario probabilities follow a Dirichlet prior:
\begin{equation}
(P(S_1), P(S_2), P(S_3)) \sim \text{Dir}(\alpha_1, \alpha_2, \alpha_3)
\label{eq:aim-dirichlet}
\end{equation}
with $\boldsymbol{\alpha} = (4, 10, 6)$, reflecting prior belief
$\mathbb{E}[P(S_2)] = 0.50$.

\textbf{Interpretation:}
\begin{itemize}[nosep]
\item $S_1$ (BigTech Disruption): AI platforms replace traditional financial advisors
\item $S_2$ (Hybrid Coexistence): AI augments but does not replace human advice
\item $S_3$ (Trust Premium): Backlash against AI increases value of human trust
\end{itemize}

\textbf{Concentration Parameter:} $\alpha_0 = \sum \alpha_k = 20$, reflecting
moderate confidence. Higher $\alpha_0$ would narrow intervals.

\textbf{MC Validation:} $P(S_2) = 0.50 \; [0.29, 0.71]$ (95\% CI).
\end{tcolorbox}


\begin{tcolorbox}[colback=gray!5!white, colframe=gray!75!black,
    title=Axiom AIM-7: Segment Behavioral Heterogeneity]
\label{ax:aim-7}

\textbf{Statement:}
The customer base partitions into $|\mathcal{G}| = 4$ behaviorally distinct
segments, each characterized by a parameter vector:
\begin{equation}
\boldsymbol{\theta}_g = (\tau_g, \lambda_g, \beta_g, \sigma_g, \eta_g, \phi_g)
\quad \forall g \in \mathcal{G}
\label{eq:aim-segment-params}
\end{equation}
with $\boldsymbol{\theta}_{g_i} \neq \boldsymbol{\theta}_{g_j}$ for $i \neq j$.

\textbf{Segments:}
\begin{itemize}[nosep]
\item $g_1$: SEG-LEGACY ($n = 1.5$M) --- Long-term UBS clients
\item $g_2$: SEG-CS-MIGRATED ($n = 1.0$M) --- Former Credit Suisse clients
\item $g_3$: SEG-DIGITAL ($n = 0.8$M) --- Digital-first clients
\item $g_4$: SEG-UHNW ($n = 0.4$M) --- Ultra-high-net-worth
\end{itemize}

\textbf{Theory Support:} MS-CJ-002 (Customer Journey Segmentation) ---
behavioral parameters, not demographics, drive segment-specific responses.
\end{tcolorbox}


\begin{tcolorbox}[colback=gray!5!white, colframe=gray!75!black,
    title=Axiom AIM-8: Multiplicative Retention]
\label{ax:aim-8}

\textbf{Statement:}
Segment retention probability is multiplicative in trust and switching cost:
\begin{equation}
R_g = \tau_g \cdot \sigma_g
\label{eq:aim-retention}
\end{equation}

\textbf{EXC-5 Justification:}
V1 (Veto): $\tau_g = 0 \Rightarrow R_g = 0$ (zero trust $\Rightarrow$ guaranteed departure,
regardless of switching cost). V1 also: $\sigma_g = 0 \Rightarrow R_g = 0$
(zero switching cost $\Rightarrow$ frictionless departure regardless of trust).

\textbf{Implication:} SEG-CS-MIGRATED ($\tau = 0.55, \sigma = 0.35$) has
$R = 0.19$ --- retention is low because \emph{both} factors are weak.
SEG-DIGITAL ($\tau = 0.70, \sigma = 0.25$) has similarly low $R = 0.17$
but for a different reason (very low switching cost despite moderate trust).

\textbf{MC Validation:}
$R_{\text{CS-MIG}} = 0.19 \; [0.08, 0.33]$;
$R_{\text{DIGITAL}} = 0.17 \; [0.08, 0.31]$.
\end{tcolorbox}


\begin{tcolorbox}[colback=gray!5!white, colframe=gray!75!black,
    title=Axiom AIM-9: Scenario-Segment Value Interaction]
\label{ax:aim-9}

\textbf{Statement:}
The expected value of segment $g$ in scenario $S_k$ is:
\begin{equation}
V_{g,k} \sim \text{Beta}(\mu_{g,k}, \sigma_{g,k})
\label{eq:aim-segment-scenario}
\end{equation}
The scenario-weighted expected value is:
\begin{equation}
\mathbb{E}[V_g] = \sum_{k} P(S_k) \cdot V_{g,k}
\label{eq:aim-expected-value}
\end{equation}

\textbf{Interpretation:} Each segment responds differently to each scenario.
SEG-DIGITAL benefits most from $S_1$ (BigTech) but suffers in $S_3$ (Trust);
SEG-UHNW benefits most from $S_3$ (Trust) but less from $S_1$.

\textbf{MC Validation:}
$\mathbb{E}[V_{\text{UHNW}}] = 0.80 \; [0.70, 0.88]$ (highest);
$\mathbb{E}[V_{\text{CS-MIG}}] = 0.54 \; [0.40, 0.69]$ (lowest).
\end{tcolorbox}


\begin{tcolorbox}[colback=gray!5!white, colframe=gray!75!black,
    title=Axiom AIM-10: Segment-Weighted Module 2 Value]
\label{ax:aim-10}

\textbf{Statement:}
Module 2 value aggregates across segments proportional to economic weight:
\begin{equation}
V_{M2} = \sum_{g \in \mathcal{G}} \frac{n_g \cdot \text{AuM}_g}{\sum_h n_h \cdot \text{AuM}_h} \cdot \mathbb{E}[V_g]
\label{eq:aim-m2-total}
\end{equation}

\textbf{Interpretation:} SEG-UHNW has disproportionate weight due to high
AuM per client, even though $n_{\text{UHNW}} < n_{\text{CS-MIG}}$.

\textbf{EXC-1 Justification:} Additivity across segments (no inter-segment
complementarity at this level; cross-module $\gamma$ captures organizational effects).
\end{tcolorbox}


%%%%%%%%%%%%%%%%%%%%%%%%%%%%%%%%%%%%%%%%%%%%%%%%%%%%%%%%%%%%%%%%%%%%%%%%%%%%%%%
\section{Module 3: Organization Transformation Axioms}
\label{sec:aim-m3}
%%%%%%%%%%%%%%%%%%%%%%%%%%%%%%%%%%%%%%%%%%%%%%%%%%%%%%%%%%%%%%%%%%%%%%%%%%%%%%%

\begin{tcolorbox}[colback=gray!5!white, colframe=gray!75!black,
    title=Axiom AIM-11: Skill Production Function]
\label{ax:aim-11}

\textbf{Statement:}
AI skill capital evolves according to a CES production function
(Cunha \& Heckman, 2007):
\begin{equation}
K_{\text{AI}}(t+1) = \left[\gamma \cdot K_{\text{AI}}(t)^{\rho} + (1-\gamma) \cdot I(t)^{\rho}\right]^{1/\rho}
\label{eq:aim-skill-production}
\end{equation}
with $\rho \in (-\infty, 1]$ governing substitutability between existing
skills and new investment.

\textbf{Theory Support:} MS-SF-001 (Self-Productivity of Skill Formation)
--- existing skills make future skill acquisition more productive.

\textbf{Implication:} Organizations with higher $K_{\text{AI}}(t_0)$ at
the start of AI transformation have a \emph{compounding advantage} in
skill acquisition.
\end{tcolorbox}


\begin{tcolorbox}[colback=gray!5!white, colframe=gray!75!black,
    title=Axiom AIM-12: Dynamic Complementarity]
\label{ax:aim-12}

\textbf{Statement:}
The marginal return on investment $I(t)$ increases with the existing skill
stock $K_{\text{AI}}(t)$:
\begin{equation}
\frac{\partial^2 K_{\text{AI}}(t+1)}{\partial I(t) \, \partial K_{\text{AI}}(t)} > 0
\label{eq:aim-dynamic-complementarity}
\end{equation}

\textbf{Corollary (DC Premium):} Investment at $t_0$ (early) yields higher
returns than the same investment at $t_0 + \Delta$ (late):
\begin{equation}
\xi = \frac{\text{ROI}(I, t_0)}{\text{ROI}(I, t_0 + \Delta)} > 1
\label{eq:aim-dc-premium}
\end{equation}

\textbf{Theory Support:} MS-SF-003 (Dynamic Complementarity) --- Heckman
(2007, 2025).

\textbf{MC Validation:} $\xi = 1.45 \; [1.26, 1.65]$;
$P(\xi > 1) = 100\%$ across 10,000 draws.
\end{tcolorbox}


\begin{tcolorbox}[colback=gray!5!white, colframe=gray!75!black,
    title=Axiom AIM-13: Role Transformation Continuity]
\label{ax:aim-13}

\textbf{Statement:}
Role transformation follows a continuous mapping, not discrete replacement:
\begin{equation}
\text{Role}(t+1) = (1 - \eta_{\text{AI}}) \cdot \text{Role}_{\text{current}}(t) + \eta_{\text{AI}} \cdot \text{Role}_{\text{augmented}}(t)
\label{eq:aim-role-transformation}
\end{equation}
where $\eta_{\text{AI}} \in [0,1]$ is the AI augmentation rate.

\textbf{Interpretation:} No marketing role disappears instantaneously.
Rather, roles evolve along a continuum from current state to AI-augmented
state. This applies to all 6 identified role transformations
(Content Creator $\to$ AI Content Curator, etc.).

\textbf{Theory Support:} MS-TP-001 (Productivity Bandwagon) --- Acemoglu (2023):
Optimal AI deployment creates new tasks alongside automating existing ones.
Pure automation without new task creation leads to ``so-so automation.''

\textbf{Anti-Pattern:} $\eta_{\text{AI}} = 1$ (complete replacement) without
new task creation $\Rightarrow$ so-so automation (Axiom AIM-14).
\end{tcolorbox}


\begin{tcolorbox}[colback=gray!5!white, colframe=gray!75!black,
    title=Axiom AIM-14: So-So Automation Avoidance]
\label{ax:aim-14}

\textbf{Statement:}
For each automated task, there must exist a complementary new task:
\begin{equation}
\forall \text{Task}_{\text{automated}}: \exists \text{Task}_{\text{new}} \text{ s.t. } V(\text{Task}_{\text{new}}) \geq V(\text{Task}_{\text{automated}})
\label{eq:aim-so-so}
\end{equation}

\textbf{Interpretation:} Automation that merely displaces workers without
creating higher-value tasks is suboptimal (``so-so automation''). The model
requires that AI deployment generates new capabilities (e.g., AI Ethics
Officer, Behavioral Intelligence Analyst) that exceed the value of
displaced activities.

\textbf{Theory Support:} MS-TP-001 (Productivity Bandwagon),
MS-TP-003 (Countervailing Powers).

\textbf{Testable Prediction:} Marketing organizations that automate without
creating new roles will show declining marginal productivity within 18 months.
\end{tcolorbox}


\begin{tcolorbox}[colback=gray!5!white, colframe=gray!75!black,
    title=Axiom AIM-15: Optimal Training Sequence]
\label{ax:aim-15}

\textbf{Statement:}
The optimal training sequence follows the dynamic complementarity hierarchy:
\begin{equation}
I^*: \quad I_{\text{literacy}} \prec I_{\text{tools}} \prec I_{\text{strategy}}
\label{eq:aim-training-sequence}
\end{equation}
where $\prec$ denotes ``must precede.'' The sequence is strict:
$\text{ROI}(I_{\text{tools}} | I_{\text{literacy}}) > \text{ROI}(I_{\text{tools}} | \neg I_{\text{literacy}})$.

\textbf{Interpretation:} AI literacy is the foundation that makes
subsequent investments in role-specific AI tools and advanced AI strategy
more productive. Investing in strategy training before literacy is suboptimal.

\textbf{Budget Recommendation:} 5--8\% of marketing budget allocated to
AI upskilling in Year 1 (Phase 1).

\textbf{Theory Support:} MS-SF-003 (Dynamic Complementarity) ---
skill $K_1$ makes acquisition of skill $K_2$ more productive.
\end{tcolorbox}


%%%%%%%%%%%%%%%%%%%%%%%%%%%%%%%%%%%%%%%%%%%%%%%%%%%%%%%%%%%%%%%%%%%%%%%%%%%%%%%
\section{Module 4: Strategic Roadmap Axioms}
\label{sec:aim-m4}
%%%%%%%%%%%%%%%%%%%%%%%%%%%%%%%%%%%%%%%%%%%%%%%%%%%%%%%%%%%%%%%%%%%%%%%%%%%%%%%

\begin{tcolorbox}[colback=gray!5!white, colframe=gray!75!black,
    title=Axiom AIM-16: No-Regret Scoring Formula]
\label{ax:aim-16}

\textbf{Statement:}
The no-regret score of a strategic move $j$ is:
\begin{equation}
\text{NR}_j = \rho_j \cdot V_j^{\min} + (1 - \rho_j) \cdot O_j + r_j \cdot L_j
\label{eq:aim-nr-score}
\end{equation}

\textbf{Interpretation:}
\begin{itemize}[nosep]
\item $\rho_j \cdot V_j^{\min}$: Guaranteed value floor (robustness $\times$ worst case)
\item $(1 - \rho_j) \cdot O_j$: Option value under uncertainty
\item $r_j \cdot L_j$: Reversible learning value
\end{itemize}

\textbf{EXC-1 Justification:} Additivity is default. The three components
are independent value sources (no veto: a move with $V_j^{\min} = 0$ can
still have $O_j > 0$).

\textbf{Design Principle:} The formula deliberately over-weights guaranteed
value and reversible learning relative to pure option value, reflecting
the behavioral insight that organizations under-explore due to
loss aversion ($\lambda > 1$).
\end{tcolorbox}


\begin{tcolorbox}[colback=gray!5!white, colframe=gray!75!black,
    title=Axiom AIM-17: No-Regret Threshold]
\label{ax:aim-17}

\textbf{Statement:}
A move qualifies as ``no-regret'' if and only if:
\begin{equation}
\text{NR}_j \geq \bar{T}_{\text{NR}} \quad \text{with} \quad P(\text{NR}_j \geq \bar{T}_{\text{NR}}) \geq 0.95
\label{eq:aim-nr-threshold}
\end{equation}
where $\bar{T}_{\text{NR}} = 0.80$ is the minimum acceptable score.

\textbf{MC Validation:} All 6 no-regret moves satisfy $P(\text{NR}_j > 0.80) = 100\%$
across 10,000 draws. The classification is robust.

\textbf{Testable Prediction:} No move with $\text{NR}_j < 0.80$ should be
initiated before scenario resolution.
\end{tcolorbox}


\begin{tcolorbox}[colback=gray!5!white, colframe=gray!75!black,
    title=Axiom AIM-18: Early Warning Signal Model]
\label{ax:aim-18}

\textbf{Statement:}
Each trigger signal $T_i(t)$ follows a stochastic growth process:
\begin{equation}
T_i(t) = T_i(0) + \mu_i \cdot t + \sigma_i \cdot W(t)
\label{eq:aim-trigger-process}
\end{equation}
where $\mu_i$ is the drift rate, $\sigma_i$ is volatility, and $W(t)$ is
standard Brownian motion.

\textbf{Scenario Detection:}
Scenario $S_k$ is ``active'' when $\geq 2$ triggers cross their $S_k$-thresholds:
\begin{equation}
\text{Active}(S_k, t) := \left|\{i : T_i(t) \geq \bar{T}_i^{(k)}\}\right| \geq 2
\label{eq:aim-scenario-detection}
\end{equation}

\textbf{Implication:} A single trigger crossing is insufficient for scenario
activation (reduces false positives). Two correlated crossings provide
statistical confidence.
\end{tcolorbox}


\begin{tcolorbox}[colback=gray!5!white, colframe=gray!75!black,
    title=Axiom AIM-19: Implementation Sequencing Optimality]
\label{ax:aim-19}

\textbf{Statement:}
The optimal implementation sequence maximizes cumulative learning value
subject to resource constraints:
\begin{equation}
\max_{\pi} \sum_{t=1}^{T} \beta^{t-1} \cdot \text{NR}_{\pi(t)} \cdot (1 + \gamma_{\text{learn}} \cdot C_{\pi}(t))
\label{eq:aim-sequencing}
\end{equation}
where $\pi$ is a permutation of moves, $\beta$ is the discount factor,
and $C_{\pi}(t) = \sum_{s<t} L_{\pi(s)}$ is cumulative learning from
prior moves.

\textbf{Implication:} Moves with high $L_j$ (learning value) should be
scheduled early to maximize the learning premium $\gamma_{\text{learn}} \cdot C$.
This yields the phase structure:
\begin{align}
\text{Phase 1:} & \quad \{NR1, NR2, NR5\} \quad \text{(high $L$, low cost)} \label{eq:aim-phase1} \\
\text{Phase 2:} & \quad \{NR3, NR6\} \quad \text{(leverages Phase 1 learning)} \label{eq:aim-phase2} \\
\text{Phase 3:} & \quad \{NR4\} \quad \text{(requires Phase 1-2 learnings)} \label{eq:aim-phase3}
\end{align}
\end{tcolorbox}


\begin{tcolorbox}[colback=gray!5!white, colframe=gray!75!black,
    title=Axiom AIM-20: Approval Threshold for Scaling]
\label{ax:aim-20}

\textbf{Statement:}
A pilot initiative transitions to full-scale deployment if and only if:
\begin{equation}
P(\text{Success} \,|\, \text{Pilot Data}) \geq \bar{A} \in [0.60, 0.80]
\label{eq:aim-approval}
\end{equation}

\textbf{Theory Support:} MS-ENT-002 (Product-Market Fit) --- Zellweger (2023):
The approval threshold reflects Bayesian updating of prior beliefs about
initiative success. The range $[0.60, 0.80]$ balances Type I (premature scaling)
and Type II (missed opportunity) errors.

\textbf{Application to NR3:} AI-Mediated Discovery Pilot requires 500 test
queries with $\bar{A} = 0.65$ (UBS mention rate in AI financial advisor responses).
\end{tcolorbox}


%%%%%%%%%%%%%%%%%%%%%%%%%%%%%%%%%%%%%%%%%%%%%%%%%%%%%%%%%%%%%%%%%%%%%%%%%%%%%%%
\section{Cross-Module Complementarity Axioms}
\label{sec:aim-cross}
%%%%%%%%%%%%%%%%%%%%%%%%%%%%%%%%%%%%%%%%%%%%%%%%%%%%%%%%%%%%%%%%%%%%%%%%%%%%%%%

\begin{definition}[Cross-Module Complementarity]
\label{def:aim-cross-module}
For modules $M_i$ and $M_j$ with values $V_i, V_j$, the complementarity
contribution is:
\begin{equation}
C_{ij} = \gamma_{ij} \cdot \sqrt{V_i \cdot V_j}
\label{eq:aim-complementarity}
\end{equation}
where $\gamma_{ij} \in [-1, 1]$. Positive $\gamma$ indicates synergy;
negative $\gamma$ indicates crowding-out.
\end{definition}

\begin{tcolorbox}[colback=red!5!white, colframe=red!75!black,
    title=Fundamental Question: Why $\sqrt{V_i \cdot V_j}$ and not $V_i \cdot V_j$?]

The geometric mean $\sqrt{V_i \cdot V_j}$ is used rather than the product
$V_i \cdot V_j$ because:
\begin{enumerate}[nosep]
\item \textbf{Scale invariance:} $\sqrt{V_i \cdot V_j}$ has the same units as $V_i$ and $V_j$
\item \textbf{Bounded contribution:} $\sqrt{V_i \cdot V_j} \leq \frac{V_i + V_j}{2}$ (AM-GM inequality)
\item \textbf{Consistency with EBF:} Appendix B uses $\gamma_{ij} \cdot \sqrt{u_i \cdot u_j}$
for utility complementarity
\end{enumerate}
\end{tcolorbox}


\begin{table}[htbp]
\centering
\caption{Cross-Module Complementarity Matrix with EXC-5 Justification}
\label{tab:aim-gamma-matrix}
\renewcommand{\arraystretch}{1.3}
\begin{tabular}{@{}llcll@{}}
\toprule
\textbf{Pair} & $\boldsymbol{\gamma}$ \textbf{[95\% CI]} & \textbf{EXC-5} & \textbf{Mechanism} & \textbf{PAR Ref} \\
\midrule
$M1 \times M3$ & $0.45 \; [0.26, 0.64]$ & V1 & Skills=0 $\Rightarrow$ AI value=0 & PAR-COMP-001 \\
$M1 \times M2$ & $0.35 \; [0.19, 0.51]$ & C1 & Customer readiness gates AI & PAR-COMP-001 \\
$M2 \times M4$ & $0.30 \; [0.14, 0.46]$ & C1 & Monitoring enables strategy & MS-ENT-001 \\
$M3 \times M4$ & $0.25 \; [0.11, 0.39]$ & S1 & Readiness scales execution & MS-SF-003 \\
\bottomrule
\end{tabular}
\end{table}


%%%%%%%%%%%%%%%%%%%%%%%%%%%%%%%%%%%%%%%%%%%%%%%%%%%%%%%%%%%%%%%%%%%%%%%%%%%%%%%
\section{Integrated Model}
\label{sec:aim-integrated}
%%%%%%%%%%%%%%%%%%%%%%%%%%%%%%%%%%%%%%%%%%%%%%%%%%%%%%%%%%%%%%%%%%%%%%%%%%%%%%%

\begin{definition}[Integrated Portfolio Value]
\label{def:aim-portfolio}
The total model value under scenario $S_k$ is:
\begin{equation}
U(S_k) = \underbrace{\sum_{m=1}^{4} w_m(S_k) \cdot V_m}_{\text{additive (EXC-1)}}
+ \underbrace{\sum_{i < j} \gamma_{ij} \cdot \sqrt{V_i \cdot V_j}}_{\text{complementarity (EXC-5)}}
\label{eq:aim-total-value}
\end{equation}
The expected portfolio value is:
\begin{equation}
U_{\text{total}} = \sum_{k=1}^{3} P(S_k) \cdot U(S_k)
\label{eq:aim-expected-total}
\end{equation}
\end{definition}

\begin{table}[htbp]
\centering
\caption{Scenario-Dependent Module Weights}
\label{tab:aim-weights}
\renewcommand{\arraystretch}{1.3}
\begin{tabular}{@{}lcccc@{}}
\toprule
\textbf{Scenario} & $w_{M1}$ & $w_{M2}$ & $w_{M3}$ & $w_{M4}$ \\
\midrule
$S_1$ BigTech Disruption & 0.35 & 0.25 & 0.25 & 0.15 \\
$S_2$ Hybrid Coexistence & 0.25 & 0.25 & 0.25 & 0.25 \\
$S_3$ Trust Premium & 0.15 & 0.20 & 0.35 & 0.30 \\
\bottomrule
\end{tabular}
\end{table}

\textbf{Weight Rationale:}
\begin{itemize}[nosep]
\item $S_1$: AI value chain ($M1$) dominates --- technology is the differentiator
\item $S_2$: Equal weights --- balanced transformation across all dimensions
\item $S_3$: Organization ($M3$) and Roadmap ($M4$) dominate --- human capital and strategic patience are key
\end{itemize}


%%%%%%%%%%%%%%%%%%%%%%%%%%%%%%%%%%%%%%%%%%%%%%%%%%%%%%%%%%%%%%%%%%%%%%%%%%%%%%%
\section{Propositions (Derived from Axiom System)}
\label{sec:aim-propositions}
%%%%%%%%%%%%%%%%%%%%%%%%%%%%%%%%%%%%%%%%%%%%%%%%%%%%%%%%%%%%%%%%%%%%%%%%%%%%%%%

\begin{proposition}[Must-Win Functions]
\label{prop:aim-must-win}
From AIM-1, AIM-3, and AIM-4:
\begin{equation}
V_{F_3} > V_{F_5} \gg V_{F_2}
\quad \text{with} \quad P(V_{F_3} \text{ in Top 3}) \geq 0.99
\end{equation}

\textbf{Proof:} By AIM-3, $\alpha_{F_3} > \alpha_{F_5}$. By AIM-1, value is
monotone in $\alpha$ (given comparable $\varphi, \delta, w$). By AIM-4,
$\delta_{F_5} > \delta_{F_3}$, which partially compensates but does not
reverse the ordering because $\alpha_{F_3} - \alpha_{F_5} > \delta_{F_5} - \delta_{F_3}$
in the parameter region.

\textbf{MC Confirmation:} $F_3$ ranked \#1 in 99.8\% of draws. \qed
\end{proposition}


\begin{proposition}[CS-Migrated Critical Risk]
\label{prop:aim-cs-risk}
From AIM-7 and AIM-8:
\begin{equation}
R_{g_2} = \min_{g \in \mathcal{G}} \{R_g\}
\quad \Leftrightarrow \quad
\text{SEG-CS-MIGRATED has highest flight risk}
\end{equation}

\textbf{Proof:} By AIM-8, $R_g = \tau_g \cdot \sigma_g$.
$R_{g_2} = 0.55 \times 0.35 = 0.19$. For all other segments,
$\tau \cdot \sigma > 0.19$ (SEG-LEGACY: 0.64, SEG-UHNW: 0.64,
SEG-DIGITAL: 0.17). Note: SEG-DIGITAL has comparably low $R$ (0.17)
but from different parametric composition.

\textbf{MC Confirmation:} Flight risk SEG-CS-MIG $= 0.81 \; [0.67, 0.92]$. \qed
\end{proposition}


\begin{proposition}[Early Training Dominance]
\label{prop:aim-early-training}
From AIM-11 and AIM-12:
\begin{equation}
P(\xi > 1) = 1.0
\quad \text{i.e., early training always dominates late training}
\end{equation}

\textbf{Proof:} By AIM-12, $\frac{\partial^2 K}{\partial I \, \partial K} > 0$.
Since $K(t_0) < K(t_0 + \Delta)$ (skills accumulate), and by AIM-11 the
production function is CES with complementarity ($\rho < 1$), the marginal
return to $I$ is increasing in the base $K$. Early investment builds a
higher $K$ for subsequent periods, creating a compounding advantage.

\textbf{MC Confirmation:} $\xi = 1.45 \; [1.26, 1.65]$;
$P(\xi > 1) = 100\%$ across all draws. \qed
\end{proposition}


\begin{proposition}[NR1 and NR2 Priority]
\label{prop:aim-nr-priority}
From AIM-16, AIM-17, and AIM-19:
\begin{equation}
\text{NR}_1 \geq \text{NR}_j \quad \forall j \in \{2,\ldots,6\}
\quad \text{with} \quad P(\text{NR}_1 = \text{rank } 1) \geq 0.60
\end{equation}

\textbf{Proof:} By AIM-16, $\text{NR}_1$ has the highest $\rho \cdot V^{\min}$
(0.95 $\times$ 0.80 = 0.76) and the highest $r \cdot L$ (0.90 $\times$ 0.90 = 0.81).
By AIM-19, high $L_1$ also generates the highest learning premium for
subsequent moves.

\textbf{MC Confirmation:} NR1 ranked \#1 in 60.8\%, NR2 in 38.8\%.
Together they cover 99.6\% of \#1 rankings. \qed
\end{proposition}


\begin{proposition}[Strongest Complementarity: Value Chain $\times$ Organization]
\label{prop:aim-strongest-gamma}
From the cross-module axiom system:
\begin{equation}
\gamma_{M1,M3} = \max_{i<j} \gamma_{ij} = 0.45 \; [0.26, 0.64]
\end{equation}

\textbf{Interpretation:} The strongest synergy is between AI capability
investment (M1) and organizational skill readiness (M3). This follows
from the veto logic: an organization without AI skills ($K_{\text{AI}} = 0$)
cannot extract value from AI capabilities ($V_{M1}$), regardless of how
advanced those capabilities are.

\textbf{Implication:} Phase 1 must address both M1 (via NR3 Discovery Pilot)
\emph{and} M3 (via NR1 AI Literacy) simultaneously.
\end{proposition}


\begin{proposition}[Scenario Sensitivity Dominance]
\label{prop:aim-sensitivity}
From AIM-6 and the integrated model:
\begin{equation}
\left|\text{Corr}(P(S_1), U_{\text{total}})\right| > \left|\text{Corr}(P(S_k), U_{\text{total}})\right|
\quad \forall k \neq 1
\end{equation}

\textbf{Interpretation:} The BigTech Disruption scenario ($S_1$) is the
strongest negative driver of total portfolio value. This is because $S_1$
shifts weight toward $M1$ (which has highest uncertainty) and away from
$M3$ (which provides stable value through human capital).

\textbf{MC Confirmation:} $r(P(S_1), U) = -0.526$;
$r(P(S_3), U) = +0.332$. \qed
\end{proposition}


\begin{proposition}[Early Warning Detection]
\label{prop:aim-early-warning}
From AIM-18:
\begin{equation}
P(\text{Active}(S_2, 24\text{m})) > 0.90
\quad \text{for triggers } T_1, T_3
\end{equation}

\textbf{Interpretation:} The Hybrid Coexistence scenario ($S_2$) will likely
be detectable within 24 months via T1 (AI Query Share) and T3 (Client AI
Adoption). This means the monitoring system (NR2) must be operational
within 6 months to capture the signal.

\textbf{MC Confirmation:}
$P(T_1 \geq \bar{T}_1^{(S_2)}, 24\text{m}) = 96.2\%$;
$P(T_3 \geq \bar{T}_3^{(S_2)}, 24\text{m}) = 92.7\%$. \qed
\end{proposition}


\begin{proposition}[Complementarity Premium Magnitude]
\label{prop:aim-comp-premium}
From the integrated model (Definition~\ref{def:aim-portfolio}):
\begin{equation}
\frac{U_{\text{total}}^{\text{with } \gamma}}{U_{\text{total}}^{\gamma = 0}} \approx 1.28 / 1.00 = 1.28
\end{equation}

\textbf{Interpretation:} Cross-module complementarities contribute approximately
28\% of total portfolio value. Ignoring complementarities (treating modules
as independent) underestimates the model value by approximately one quarter.

\textbf{Implication:} Sequential, siloed implementation of modules
(e.g., ``first M1, then M3'') destroys complementarity value.
Parallel implementation preserves $\gamma_{ij}$ contributions.
\end{proposition}


%%%%%%%%%%%%%%%%%%%%%%%%%%%%%%%%%%%%%%%%%%%%%%%%%%%%%%%%%%%%%%%%%%%%%%%%%%%%%%%
\section{Axiom Dependency Graph}
\label{sec:aim-dependency}
%%%%%%%%%%%%%%%%%%%%%%%%%%%%%%%%%%%%%%%%%%%%%%%%%%%%%%%%%%%%%%%%%%%%%%%%%%%%%%%

\begin{figure}[htbp]
\centering
\begin{verbatim}
    AIM-1 ──┬── AIM-3 ──── Prop 1 (Must-Win Functions)
            │
    AIM-2 ──┘
            ├── AIM-4
            │
    AIM-5 ──┘

    AIM-6 ──┬── AIM-9 ──── Prop 6 (Sensitivity Dominance)
            │
    AIM-7 ──┼── AIM-8 ──── Prop 2 (CS-Migrated Risk)
            │
    AIM-10 ─┘

    AIM-11 ─┬── AIM-12 ──── Prop 3 (Early Training)
            │
    AIM-13 ─┼── AIM-14
            │
    AIM-15 ─┘

    AIM-16 ─┬── AIM-17 ──── Prop 4 (NR Priority)
            │
    AIM-18 ─┼─────────────── Prop 7 (Early Warning)
            │
    AIM-19 ─┼── AIM-20
            │
    gamma ──┘── Prop 5 (Strongest gamma)
                Prop 8 (Complementarity Premium)
\end{verbatim}
\caption{Axiom Dependency Graph for MOD-UBS-AIMARK-001}
\label{fig:aim-dependency}
\end{figure}


%%%%%%%%%%%%%%%%%%%%%%%%%%%%%%%%%%%%%%%%%%%%%%%%%%%%%%%%%%%%%%%%%%%%%%%%%%%%%%%
\section{Summary}
\label{sec:aim-summary}
%%%%%%%%%%%%%%%%%%%%%%%%%%%%%%%%%%%%%%%%%%%%%%%%%%%%%%%%%%%%%%%%%%%%%%%%%%%%%%%

\begin{tcolorbox}[colback=blue!5!white, colframe=blue!75!black,
    title=Summary: AI Marketing Transformation Axiom System]

\textbf{20 Axioms} organized in 4 modules + cross-module:
\begin{itemize}[nosep]
\item \textbf{M1} (AIM-1 to AIM-5): Value decomposition, function space, AI ordering
\item \textbf{M2} (AIM-6 to AIM-10): Scenarios, segments, retention, interaction
\item \textbf{M3} (AIM-11 to AIM-15): Skill formation, DC, role transformation
\item \textbf{M4} (AIM-16 to AIM-20): NR scoring, early warning, sequencing
\end{itemize}

\textbf{8 Propositions} all validated via Monte Carlo (10,000 draws):
\begin{enumerate}[nosep]
\item $F_3$ is must-win ($P = 99.8\%$ top 3)
\item SEG-CS-MIGRATED is highest flight risk ($R = 0.19$)
\item Early training always dominates ($P(\xi > 1) = 100\%$)
\item NR1 + NR2 are priority moves ($99.6\%$ rank \#1)
\item $\gamma_{M1,M3} = 0.45$ is the strongest complementarity
\item $P(S_1)$ is the dominant risk driver ($r = -0.526$)
\item $S_2$ detectable within 24 months ($P > 90\%$)
\item Complementarity premium $\approx 28\%$ of total value
\end{enumerate}

\textbf{Key Principle:} All propositions are \emph{derived} from the axiom system
and \emph{confirmed} by Monte Carlo simulation. The axiom system is
falsifiable: empirical data from UBS pilots (Q3 2026) can test each
proposition independently.
\end{tcolorbox}


%%%%%%%%%%%%%%%%%%%%%%%%%%%%%%%%%%%%%%%%%%%%%%%%%%%%%%%%%%%%%%%%%%%%%%%%%%%%%%%
\section{Glossary of Symbols}
\label{sec:aim-glossary}
%%%%%%%%%%%%%%%%%%%%%%%%%%%%%%%%%%%%%%%%%%%%%%%%%%%%%%%%%%%%%%%%%%%%%%%%%%%%%%%

See Appendix~\ref{app:glossary} (G) for the complete EBF glossary.
Table~\ref{tab:aim-symbols} provides all symbols specific to this appendix.


%%%%%%%%%%%%%%%%%%%%%%%%%%%%%%%%%%%%%%%%%%%%%%%%%%%%%%%%%%%%%%%%%%%%%%%%%%%%%%%
\section{Open Issues and Research Agenda}
\label{sec:aim-open}
%%%%%%%%%%%%%%%%%%%%%%%%%%%%%%%%%%%%%%%%%%%%%%%%%%%%%%%%%%%%%%%%%%%%%%%%%%%%%%%

\begin{enumerate}
\item \textbf{Empirical calibration:} AIM-8 (multiplicative retention) is an
assumption that requires validation with longitudinal client data.
Alternative: additive $R_g = a \cdot \tau_g + b \cdot \sigma_g$.

\item \textbf{Intra-module complementarity:} AIM-5 assumes additivity within M1.
If $F_3$ and $F_5$ exhibit synergy (AI Intelligence improves AI Discovery),
an intra-module $\gamma_{F_3, F_5}$ should be added.

\item \textbf{Scenario transition:} AIM-6 treats scenarios as static.
A dynamic model with transition probabilities $P(S_k(t+1) | S_j(t))$
would capture scenario evolution.

\item \textbf{Cross-segment effects:} AIM-10 assumes segment independence.
If SEG-DIGITAL defection creates negative peer effects for SEG-LEGACY
(social proof of switching), a cross-segment $\gamma$ is needed.

\item \textbf{Adversarial dynamics:} The model does not capture competitor
responses. A game-theoretic extension (BigTech vs.\ UBS in $S_1$) would
improve M4 predictions.
\end{enumerate}


%%%%%%%%%%%%%%%%%%%%%%%%%%%%%%%%%%%%%%%%%%%%%%%%%%%%%%%%%%%%%%%%%%%%%%%%%%%%%%%
\section{References}
\label{sec:aim-references}
%%%%%%%%%%%%%%%%%%%%%%%%%%%%%%%%%%%%%%%%%%%%%%%%%%%%%%%%%%%%%%%%%%%%%%%%%%%%%%%

\textbf{Master Bibliography:} See \nocite{bcm_master} (\textit{bcm\_master.bib})

\textbf{Key References:}
\begin{itemize}[nosep]
\item \citet{acemoglu2023power} --- Productivity Bandwagon, so-so automation (AIM-13, AIM-14)
\item \citet{cunha2007technology} --- Skill Formation Technology (AIM-11, AIM-12)
\item \citet{heckman2025dynamic} --- Dynamic Complementarity (AIM-12, AIM-15)
\item \citet{puntoni2021consumers} --- AI Consumer Experience (AIM-3, AIM-4)
\item \citet{akerlof2000economics} --- Identity Economics, segment identity (AIM-7)
\item \citet{zellweger2023pragmatist} --- Pragmatist Entrepreneurship, approval threshold (AIM-20)
\item \citet{fehr1999theory} --- Inequity Aversion, trust dynamics (AIM-8)
\item \citet{herhausen2019loyalty} --- Digital loyalty, switching behavior (AIM-8)
\end{itemize}


%%%%%%%%%%%%%%%%%%%%%%%%%%%%%%%%%%%%%%%%%%%%%%%%%%%%%%%%%%%%%%%%%%%%%%%%%%%%%%%
\section{Cross-Reference Footer}
\label{sec:aim-xref}
%%%%%%%%%%%%%%%%%%%%%%%%%%%%%%%%%%%%%%%%%%%%%%%%%%%%%%%%%%%%%%%%%%%%%%%%%%%%%%%

This appendix is referenced by:
\begin{itemize}
\item Model Registry: MOD-UBS-AIMARK-001 (v1.1)
\item Session: EBF-S-2026-02-12-FIN-001
\item Output Registry: EBF-OUT-FIN-001
\item Monte Carlo Script: \texttt{scripts/monte\_carlo\_ubs\_aimark.py}
\end{itemize}

This appendix references:
\begin{itemize}
\item Appendix B (CORE-HOW): General complementarity framework ($\gamma_{ij}$)
\item Appendix V (CORE-WHEN): Context dimensions ($\Psi$)
\item Appendix BBB (CORE-WHERE): Parameter sources
\item Appendix AV (CORE-READY): Skill readiness framework
\item Appendix IE (CORE-EIT): Intervention emergence methodology
\item Appendix SF: Dynamic Complementarity (Cunha/Heckman)
\end{itemize}
